% !PRINTED 341
\weberpage{341}
% running head: §. 91. | Realität. | 341
%<<CONT

\weq{18}{\begin{array}{ccc}
(0, 1, -1), & (-\varepsilon, 0, 1), & (1, -\varepsilon^2, 0) \\
(-1, 0, 1), & (1, -\varepsilon, 0), & (0, 1, -\varepsilon^2) \\
(1, -1, 0), & (0, 1, -\varepsilon), & (-\varepsilon^2, 0, 1),
\end{array}}
\weq{19}{\begin{array}{ccc}
(0, 1, -1), & (-\varepsilon^2, 0, 1), & (1, -\varepsilon, 0) \\
(-1, 0, 1), & (1, -\varepsilon^2, 0), & (0, 1, -\varepsilon) \\
(1, -1, 0), & (0, 1, -\varepsilon^2), & (-\varepsilon, 0, 1).
\end{array}}

Von den geraden Linien (16) enthält jede einen reellen
und zwei conjugirt imaginäre Punkte, und alle drei Linien sind
reell. Von den Geraden (17) enthält die erste die drei reellen
Punkte und ist reell, während die beiden anderen conjugirt
imaginär sind. In (18) und (19) sind alle drei Linien imaginär,
und zwar sind die Linien von (18) conjugirt zu den Linien
von (19).

Eine noch übersichtlichere Bezeichnung ist folgende. Wenn
$\xi$ und $\eta$ je ein volles Restsystem nach dem Modul 3 durchlaufen
, etwa $0, 1, 2$, und zwei nach dem Modul 3 congruente
Zahlen als nicht verschieden angesehen werden, so giebt es neun
Combinationen $(\xi, \eta)$, die wir als Bezeichnung für die neun
Wendepunkte benutzen können. Setzen wir fest, dass $(\xi, \eta)$ und
$(\eta, \xi)$ conjugirt imaginäre Wendepunkte, also $(0, 0)$, $(1, 1)$, $(2, 2)$
die drei reellen bedeuten, so erhalten wir aus (16) bis (19) folgende
 Tabelle, in der wieder die in einer Zeile stehenden Wendepunkte
 auf einer geraden Linie liegen:
\weq{20}{\begin{array}{ccc@{\qquad}ccc}
(0,0) & (1,2) & (2,1) & (0,0) & (1,1) & (2,2) \\
(1,1) & (2,0) & (0,2) & (1,2) & (2,0) & (0,1) \\
(2,2) & (0,1) & (1,0) & (2,1) & (0,2) & (1,0) \\[4pt]
(0,0) & (2,0) & (1,0) & (0,0) & (0,2) & (0,1) \\
(1,1) & (0,1) & (2,1) & (1,1) & (1,0) & (1,2) \\
(2,2) & (1,2) & (0,2) & (2,2) & (2,1) & (2,0).
\end{array}}

Man sieht hieraus, dass drei Wendepunkte $(\xi_1, \eta_1)$, $(\xi_2, \eta_2)$,
$(\xi_3, \eta_3)$ dann und nur dann auf einer geraden Linie liegen, wenn
die beiden Congruenzen
\weq{21}{\begin{aligned}
\xi_1+\xi_2+\xi_3 &\equiv 0 \\
\eta_1+\eta_2+\eta_3 &\equiv 0
\end{aligned} \quad (\mathrm{mod}\ 3)}
erfüllt sind.
